% Options for packages loaded elsewhere
\PassOptionsToPackage{unicode}{hyperref}
\PassOptionsToPackage{hyphens}{url}
%
\documentclass[
  17pt,
  letterpaper,
  ignorenonframetext,
  aspectratio=169,
  handout]{beamer}
\usepackage{pgfpages}
\setbeamertemplate{caption}[numbered]
\setbeamertemplate{caption label separator}{: }
\setbeamercolor{caption name}{fg=normal text.fg}
\beamertemplatenavigationsymbolsempty
% Prevent slide breaks in the middle of a paragraph
\widowpenalties 1 10000
\raggedbottom
\setbeamertemplate{part page}{
  \centering
  \begin{beamercolorbox}[sep=16pt,center]{part title}
    \usebeamerfont{part title}\insertpart\par
  \end{beamercolorbox}
}
\setbeamertemplate{section page}{
  \centering
  \begin{beamercolorbox}[sep=12pt,center]{part title}
    \usebeamerfont{section title}\insertsection\par
  \end{beamercolorbox}
}
\setbeamertemplate{subsection page}{
  \centering
  \begin{beamercolorbox}[sep=8pt,center]{part title}
    \usebeamerfont{subsection title}\insertsubsection\par
  \end{beamercolorbox}
}
\AtBeginPart{
  \frame{\partpage}
}
\AtBeginSection{
  \ifbibliography
  \else
    \frame{\sectionpage}
  \fi
}
\AtBeginSubsection{
  \frame{\subsectionpage}
}

\usepackage{amsmath,amssymb}
\usepackage{lmodern}
\usepackage{iftex}
\ifPDFTeX
  \usepackage[T1]{fontenc}
  \usepackage[utf8]{inputenc}
  \usepackage{textcomp} % provide euro and other symbols
\else % if luatex or xetex
  \usepackage{unicode-math}
  \defaultfontfeatures{Scale=MatchLowercase}
  \defaultfontfeatures[\rmfamily]{Ligatures=TeX,Scale=1}
  \setmainfont[BoldFont = SF Pro Text Semibold, Scale =
MatchLowercase]{SF Pro Text Light}
\fi
\usecolortheme{wolverine}
\usefonttheme{serif} % use mainfont rather than sansfont for slide text
\useinnertheme{default}
\useoutertheme{miniframes}
% Use upquote if available, for straight quotes in verbatim environments
\IfFileExists{upquote.sty}{\usepackage{upquote}}{}
\IfFileExists{microtype.sty}{% use microtype if available
  \usepackage[]{microtype}
  \UseMicrotypeSet[protrusion]{basicmath} % disable protrusion for tt fonts
}{}
\makeatletter
\@ifundefined{KOMAClassName}{% if non-KOMA class
  \IfFileExists{parskip.sty}{%
    \usepackage{parskip}
  }{% else
    \setlength{\parindent}{0pt}
    \setlength{\parskip}{6pt plus 2pt minus 1pt}}
}{% if KOMA class
  \KOMAoptions{parskip=half}}
\makeatother
\usepackage{xcolor}
\newif\ifbibliography
\setlength{\emergencystretch}{3em} % prevent overfull lines
\setcounter{secnumdepth}{-\maxdimen} % remove section numbering


\providecommand{\tightlist}{%
  \setlength{\itemsep}{0pt}\setlength{\parskip}{0pt}}\usepackage{longtable,booktabs,array}
\usepackage{calc} % for calculating minipage widths
\usepackage{caption}
% Make caption package work with longtable
\makeatletter
\def\fnum@table{\tablename~\thetable}
\makeatother
\usepackage{graphicx}
\makeatletter
\def\maxwidth{\ifdim\Gin@nat@width>\linewidth\linewidth\else\Gin@nat@width\fi}
\def\maxheight{\ifdim\Gin@nat@height>\textheight\textheight\else\Gin@nat@height\fi}
\makeatother
% Scale images if necessary, so that they will not overflow the page
% margins by default, and it is still possible to overwrite the defaults
% using explicit options in \includegraphics[width, height, ...]{}
\setkeys{Gin}{width=\maxwidth,height=\maxheight,keepaspectratio}
% Set default figure placement to htbp
\makeatletter
\def\fps@figure{htbp}
\makeatother

\captionsetup[figure]{labelformat=empty}
\usepackage{pgfpages}
\setbeamertemplate{itemize item}[circle]
\setbeamertemplate{footline}[frame number]{}
\mode<handout>{\pgfpagesuselayout{6 on 1}[letterpaper, border shrink=8mm]}
\AtBeginSection{%
   \begin{frame}
       \tableofcontents[currentsection]
   \end{frame}
}
\makeatletter
\makeatother
\makeatletter
\makeatother
\makeatletter
\@ifpackageloaded{caption}{}{\usepackage{caption}}
\AtBeginDocument{%
\ifdefined\contentsname
  \renewcommand*\contentsname{Table of contents}
\else
  \newcommand\contentsname{Table of contents}
\fi
\ifdefined\listfigurename
  \renewcommand*\listfigurename{List of Figures}
\else
  \newcommand\listfigurename{List of Figures}
\fi
\ifdefined\listtablename
  \renewcommand*\listtablename{List of Tables}
\else
  \newcommand\listtablename{List of Tables}
\fi
\ifdefined\figurename
  \renewcommand*\figurename{Figure}
\else
  \newcommand\figurename{Figure}
\fi
\ifdefined\tablename
  \renewcommand*\tablename{Table}
\else
  \newcommand\tablename{Table}
\fi
}
\@ifpackageloaded{float}{}{\usepackage{float}}
\floatstyle{ruled}
\@ifundefined{c@chapter}{\newfloat{codelisting}{h}{lop}}{\newfloat{codelisting}{h}{lop}[chapter]}
\floatname{codelisting}{Listing}
\newcommand*\listoflistings{\listof{codelisting}{List of Listings}}
\makeatother
\makeatletter
\@ifpackageloaded{caption}{}{\usepackage{caption}}
\@ifpackageloaded{subcaption}{}{\usepackage{subcaption}}
\makeatother
\makeatletter
\@ifpackageloaded{tcolorbox}{}{\usepackage[many]{tcolorbox}}
\makeatother
\makeatletter
\@ifundefined{shadecolor}{\definecolor{shadecolor}{rgb}{.97, .97, .97}}
\makeatother
\makeatletter
\makeatother
\ifLuaTeX
  \usepackage{selnolig}  % disable illegal ligatures
\fi
\IfFileExists{bookmark.sty}{\usepackage{bookmark}}{\usepackage{hyperref}}
\IfFileExists{xurl.sty}{\usepackage{xurl}}{} % add URL line breaks if available
\urlstyle{same} % disable monospaced font for URLs
\hypersetup{
  pdftitle={Knowledge and Reality, Lecture 06},
  pdfauthor={Brian Weatherson},
  hidelinks,
  pdfcreator={LaTeX via pandoc}}

\title{Knowledge and Reality, Lecture 06}
\author{Brian Weatherson}
\date{2022-09-19}

\begin{document}
\frame{\titlepage}
\ifdefined\Shaded\renewenvironment{Shaded}{\begin{tcolorbox}[enhanced, borderline west={3pt}{0pt}{shadecolor}, breakable, interior hidden, boxrule=0pt, sharp corners, frame hidden]}{\end{tcolorbox}}\fi

\hypertarget{review}{%
\section{Review}\label{review}}

\begin{frame}{Internalism}
\protect\hypertarget{internalism}{}
Rationality is a function of what's internal to the person.

\begin{itemize}[<+->]
\tightlist
\item
  Two people who are internally alike are either both rational, or both
  irrational.
\end{itemize}
\end{frame}

\begin{frame}{Externalism}
\protect\hypertarget{externalism}{}
Externalism is simply the denial of internalism.

\begin{itemize}[<+->]
\tightlist
\item
  It says that sometimes internal duplicates differ in rationality.
\end{itemize}
\end{frame}

\begin{frame}{Process Reliabilism}
\protect\hypertarget{process-reliabilism}{}
The most popular form of externalism these days is \textbf{process
reliabilism}. This has two parts.

\begin{itemize}[<+->]
\tightlist
\item
  In general, rational beliefs are those that are produced by reliable
  processes.
\item
  But there is an exception; beliefs that fail some internalist test
  might be `defeated', and not rational.
\end{itemize}
\end{frame}

\hypertarget{why-internalism}{%
\section{Why Internalism}\label{why-internalism}}

\begin{frame}{Three Reasons}
\protect\hypertarget{three-reasons}{}
Many reasons, we'll focus on three. The first two were mentioned last
time.

\begin{enumerate}[<+->]
\tightlist
\item
  Anti-luck.
\item
  Cases.
\item
  Evil Demons
\end{enumerate}
\end{frame}

\begin{frame}{Evil Demon}
\protect\hypertarget{evil-demon}{}
The evil demon usually comes into epistemology as part of an argument
about scepticism.

\begin{itemize}[<+->]
\tightlist
\item
  E.g., you don't know you're not being deceived by an evil demon, so
  you don't know that you have hands.
\item
  But here is gets used for a debate about rationality.
\end{itemize}
\end{frame}

\begin{frame}{Sympathy for the Devil (Victim)}
\protect\hypertarget{sympathy-for-the-devil-victim}{}
Imagine a person who seems, from the inside, to be just like an actually
rational person.

\begin{itemize}[<+->]
\tightlist
\item
  But in fact they are the victim of an evil demon, so their beliefs are
  all false.
\item
  Reliabilists say they are irrational, but they are intuitively
  rational.
\end{itemize}
\end{frame}

\begin{frame}{Question}
\protect\hypertarget{question}{}
What should reliabilists say about this case?
\end{frame}

\begin{frame}{Variant}
\protect\hypertarget{variant}{}
Imagine two evil demon victims.

\begin{itemize}[<+->]
\tightlist
\item
  A is, from the inside, a paradigm of rationality.
\item
  B is, from the inside, a conspiracy theorist who leaps to conclusions,
  believes everything he's told, etc., etc.
\end{itemize}
\end{frame}

\begin{frame}{Question}
\protect\hypertarget{question-1}{}
Can reliabilists say something about the intuitive differences between A
and B?
\end{frame}

\hypertarget{why-externalism}{%
\section{Why Externalism}\label{why-externalism}}

\begin{frame}{Three Arguments}
\protect\hypertarget{three-arguments}{}
Again, there are many, but I'll just look at these three.

\begin{enumerate}[<+->]
\tightlist
\item
  Speckled hen
\item
  Why care about rationality?
\item
  Involuntarism
\end{enumerate}
\end{frame}

\begin{frame}{Speckled Hen}
\protect\hypertarget{speckled-hen}{}
\begin{figure}

{\centering \includegraphics[width=\textwidth,height=0.6\textheight]{../images/hen.jpg}

}

\caption{How many visible white speckles?}

\end{figure}
\end{frame}

\begin{frame}{Speckled Hen}
\protect\hypertarget{speckled-hen-1}{}
\begin{itemize}[<+->]
\tightlist
\item
  The right answer does just depend on your internal states.
\item
  But it seems that any number you came up with by just guessing would
  not be a rational belief.
\item
  So rationality requires more than internal states, it requires a
  reliable connection to the world.
\end{itemize}
\end{frame}

\begin{frame}{Why Care?}
\protect\hypertarget{why-care}{}
Rationality should be something we care about.

\begin{itemize}[<+->]
\tightlist
\item
  But if given a choice between being evidence-responsive, and being
  reliable, we should choose being reliable.
\item
  And same for coherent versus reliable.
\end{itemize}
\end{frame}

\begin{frame}{Why Care?}
\protect\hypertarget{why-care-1}{}
\begin{itemize}[<+->]
\tightlist
\item
  Belief, in some sense, aims at the truth.
\item
  Rationality means something like doing well in believing.
\item
  So, being rational should mean something like doing well in getting to
  the truth.
\item
  That means reliability.
\end{itemize}
\end{frame}

\begin{frame}{Why Care?}
\protect\hypertarget{why-care-2}{}
Are there hand-wavy moves on the previous slide that might not hold up
to strict scrutiny? Yes!
\end{frame}

\begin{frame}{Voluntarism}
\protect\hypertarget{voluntarism}{}
\begin{itemize}[<+->]
\tightlist
\item
  We evaluate the voluntary parts of human behavior by whether they make
  sense, and the involuntary parts by whether they work.
\item
  Belief is involuntary.
\item
  So we should evaluate it by how well it works.
\end{itemize}
\end{frame}

\begin{frame}{Somewhat Concessive}
\protect\hypertarget{somewhat-concessive}{}
\begin{itemize}[<+->]
\tightlist
\item
  This argument concedes that if belief were voluntary, internalism
  would be plausible/correct.
\item
  But, it says, belief is involuntary.
\end{itemize}
\end{frame}

\begin{frame}{Evaluation}
\protect\hypertarget{evaluation}{}
Whether someone has, say, a good or bad digestive system is not a
strictly internal matter.

\begin{itemize}[<+->]
\tightlist
\item
  Having a good digestive system just is being good at digesting common
  foods.
\item
  And what the common foods are is external.
\item
  The involunarists say belief is the same.
\end{itemize}
\end{frame}

\begin{frame}{Two Objections}
\protect\hypertarget{two-objections}{}
\begin{enumerate}[<+->]
\tightlist
\item
  Belief really is voluntary.
\item
  Enough things connected to belief are voluntary that we can use
  something like internalist criteria.
\end{enumerate}
\end{frame}

\hypertarget{srinivasan}{%
\section{Srinivasan}\label{srinivasan}}

\begin{frame}{Short Version}
\protect\hypertarget{short-version}{}
\begin{itemize}[<+->]
\tightlist
\item
  In cases involving accurate beliefs under oppressive circumstances,
  reliability is enough for rationality.
\item
  So internalism is false as applied to these rather important actual
  cases.
\end{itemize}
\end{frame}

\begin{frame}{Three Cases}
\protect\hypertarget{three-cases}{}
\begin{enumerate}[<+->]
\tightlist
\item
  Instinctive correct belief.
\item
  Instinctive correct belief which persists despite counter-evidence.
\item
  False belief that matches lots of testimony.
\end{enumerate}
\end{frame}

\begin{frame}{Three Rival Theories}
\protect\hypertarget{three-rival-theories}{}
\begin{itemize}[<+->]
\tightlist
\item
  Internalism.
\item
  This is incompatible with all of the cases.
\end{itemize}
\end{frame}

\begin{frame}{Three Rival Theories}
\protect\hypertarget{three-rival-theories-1}{}
\begin{itemize}[<+->]
\tightlist
\item
  Process reliabilism, with the defeater condition.
\item
  This is incompatible with the second case.
\end{itemize}
\end{frame}

\begin{frame}{Three Rival Theories}
\protect\hypertarget{three-rival-theories-2}{}
\begin{itemize}[<+->]
\tightlist
\item
  I don't know of anyone who has this view, but the third case seems
  designed to oppose a view that someone probably should have.
\item
  Call it the generous view of rationality.
\end{itemize}
\end{frame}

\begin{frame}{Three Rival Theories}
\protect\hypertarget{three-rival-theories-3}{}
Here's how the `generous' view (my term) works. A belief is rational iff
it is:

\begin{enumerate}[<+->]
\tightlist
\item
  Formed by a reliable process; or
\item
  Properly based in the evidence.
\end{enumerate}
\end{frame}

\begin{frame}{Three Rival Theories}
\protect\hypertarget{three-rival-theories-4}{}
No one has this view, but it seems interesting because it gets a lot of
cases right, especially evil demon cases.

\begin{itemize}[<+->]
\tightlist
\item
  But it gets the third of Srinivasan's cases wrong.
\end{itemize}
\end{frame}

\begin{frame}{Case Judgments}
\protect\hypertarget{case-judgments}{}
What do \textbf{you} think about the three cases?
\end{frame}

\begin{frame}{Argument Structure}
\protect\hypertarget{argument-structure}{}
What role do the cases play in Srinivasan's theory?

\begin{enumerate}[<+->]
\tightlist
\item
  Trigger snap reactions that theory is judged against?
\item
  Make us reflect on what we want a theory of good belief for?
\end{enumerate}
\end{frame}

\begin{frame}{Argument Structure}
\protect\hypertarget{argument-structure-1}{}
\begin{itemize}[<+->]
\tightlist
\item
  My (idiosyncratic) view is that 1 isn't a very good reason to think
  about cases like these.
\item
  But 2 could be a good reason.
\end{itemize}
\end{frame}

\begin{frame}{Cases and Principles}
\protect\hypertarget{cases-and-principles}{}
What general lessons about belief can we draw from reflection on these
kinds of cases?
\end{frame}

\begin{frame}{Cues}
\protect\hypertarget{cues}{}
On the next slide is one kind of argument I think can be drawn from
these cases (not sure if it's a fair reading of the paper though).
\end{frame}

\begin{frame}{Cues}
\protect\hypertarget{cues-1}{}
\begin{enumerate}[<+->]
\tightlist
\item
  Rational, justified belief is a matter of doing well in believing.
\item
  Doing well in believing, for creatures like us, involves picking up on
  subtle cues.
\item
  This is often something that is not available to consciousness, or
  available to reasoning.
\item
  So rationality isn't just about reasoning from conscious states.
\end{enumerate}
\end{frame}

\begin{frame}{Racists and Clairvoyants}
\protect\hypertarget{racists-and-clairvoyants}{}
To end, let's look at Srinivasan's argument that her examples generalise
to promote a simple, or what she calls \textbf{radical} externalism.

\begin{itemize}[<+->]
\tightlist
\item
  I found this part of the argument rather odd.
\end{itemize}
\end{frame}

\begin{frame}{Argument by Analogy}
\protect\hypertarget{argument-by-analogy}{}
\begin{enumerate}[<+->]
\tightlist
\item
  Nour is justified in believing her host is racist.
\item
  Nour's case is just like Bonjour's case of Norman.
\item
  So Norman is justified in his clairvoyant beliefs.
\end{enumerate}
\end{frame}

\begin{frame}{Argument by Analogy}
\protect\hypertarget{argument-by-analogy-1}{}
\begin{itemize}[<+->]
\tightlist
\item
  Why should we believe premise 2 here?
\item
  I know why an \textbf{internalist} should believe it, by stipulation
  the cases are pretty similar from the inside, but why should an
  \textbf{externalist} believe it?
\end{itemize}
\end{frame}

\begin{frame}{Disanalogies}
\protect\hypertarget{disanalogies}{}
\begin{itemize}[<+->]
\tightlist
\item
  The Nour example is pretty realistic, the Norman case is totally not.
\item
  There is an explanation for why Nour could have this ability; there is
  no explanation for Norman.
\item
  Nour's ability is widespread; Norman's is idiosyncratic.
\end{itemize}
\end{frame}

\begin{frame}{Group Externalism}
\protect\hypertarget{group-externalism}{}
So here's a view that is untouched by Srinivasan's example, but agrees
with Bonjour about Norman.

\begin{itemize}[<+->]
\tightlist
\item
  A belief is justified iff it is produced by a process that is
  reliable, and widely shared among similar people.
\end{itemize}
\end{frame}

\begin{frame}{Group Externalism}
\protect\hypertarget{group-externalism-1}{}
This could obviously do with some more care about the details - I
literally just made it up - but it feels more natural given the Marxist
motivations to think about groups than individual reliability.
\end{frame}

\begin{frame}{For Next Time}
\protect\hypertarget{for-next-time}{}
Onto the most notorious part of recent epistemology: The Gettier Problem
\end{frame}



\end{document}
